\input{../tex-inputs/latex-leseansicht-vorspann}

\section[Gustav Schwarzkopf an Arthur Schnitzler, 5. 7. 1915]{L04246 Gustav Schwarzkopf an Arthur Schnitzler, 5. 7. 1915}
\nopagebreak\mylabel{L04246v}
\rehead{ }\normalsize\beginnumbering\briefempfaengerindex{Schnitzler, Arthur@\textsc{Schnitzler, Arthur}!zzzSchwarzkopf, Gustav@\emph{von Gustav Schwarzkopf}!1915-07-051@{5. 7. 1915}|(be}
\toendnotes[C]{\smallbreak\pagebreak[2]}
\correspDesc{Versand  durch Gustav Schwarzkopf am 5. 7. 1915 in Bad Ischl
\newline{}Erhalt  durch Arthur Schnitzler im Zeitraum [6. 7. 1915
                  – 10. 7. 1915?] in Wien}\toendnotes[C]{\smallbreak}
\Standort{CUL, Schnitzler, B 96a.}
\physDesc{Bildpostkarte, 492 Zeichen
\newline{}Handschrift: blaue Tinte, deutsche Kurrent
\newline{}Versand: Stempel: »\nobreak{}\oindex{Bad Ischl@\textbf{Bad Ischl}|pwk}Bad Ischl, 5. VII. 15, 7\nobreak{}«.  }\toendnotes[C]{\smallbreak}\pstart{}{\pb}Herrn\pend{}\pstart{}\textsc{D\textsuperscript{r} Arthur Schnitzler}\pend{}\pstart{}\textsc{Wien XVIII\oindex{XVIII., Währing@\textbf{XVIII., Währing}, \emph{Verwaltungsgebiet}|pw}}\pend{}\pstart{}Sternwarteſtraße 71\oindex{Wien@\textbf{Wien}!XVIII., Währing@\textbf{XVIII., Währing}!Sternwartestraße 71@\textbf{Sternwartestraße 71}, \emph{Wohngebäude}|pw}.\pend{}{\bigskip}
\pstart
           \noindent{}{\pb}\textcolor{gray}{\textbf{LAUFEN bei Bad Ischl\oindex{Laufen@\textbf{Laufen}|pw}, 479 m Seehöhe. Salzkammergut\oindex{Salzkammergut@\textbf{Salzkammergut}, \emph{Region}|pw}.}}\pend
           \vspace{1em}
\pstart
           9. VII. 15.\hfill \textsc{Hôtel garni Athen\oindex{Hotel garni Athen@\textbf{Hotel garni Athen}, \emph{Hotel}|pw}}, Z. 12.\pend
           \vspace{0.5em}
\pstart
           Lieber Arthur, es iſt jetzt kein Zweifel mehr erlaubt, ich bin doch
               hier. In einem{ }ſehr angenehmen Zi{\geminationm}er mit wirklich
               hübſcher Ausſicht. Nach den bisherigen Erfahrungen ka{\geminationn}
               ich B. H.\pwindex{Beer-Hofmann, Richard 11.\,7.\,1866 Wien – 26.\,9.\,1945 New York City@\textsc{Beer-Hofmann, Richard} (11.\,7.\,1866 Wien – 26.\,9.\,1945 New York City), \emph{Schriftsteller}|pw} Wahl nur loben. (Seine Wohnung\oindex{Villa Gisela@\textbf{Villa Gisela}, \emph{Wohngebäude}|pwv} liegt{ }ſehr hübſch.)
               Es iſt hier noch ſehr ſtill. Auf der Esplanade ka{\geminationn} man
               beinahe träumen, aber auch von dieſer Freiheit machen wir wenigen Gebrauch.\pend
           
\pstart
           Ihre Frau\pwindex{Schnitzler, Olga 17.\,1.\,1882 Wien – 13.\,1.\,1970 Lugano@\textsc{Schnitzler, Olga} (17.\,1.\,1882 Wien – 13.\,1.\,1970 Lugano), \emph{Schauspielerin, Sängerin}|pwv} und Sie
               herzlichſt grüßend,{\\[\baselineskip]}Ihr{\\[\baselineskip]}\spacefill\mbox{Gustav.}\pend
           \leftskip=0em{}\selectlanguage{ngerman}\endnumbering\briefempfaengerindex{Schnitzler, Arthur@\textsc{Schnitzler, Arthur}!zzzSchwarzkopf, Gustav@\emph{von Gustav Schwarzkopf}!1915-07-051@{5. 7. 1915}|)be}\mylabel{L04246h}
\begin{anhang}
\end{anhang}\newcommand{\dateiname}{L04246}\newcommand{\titel}{Gustav Schwarzkopf an Arthur Schnitzler, 5. 7. 1915}\newcommand{\editorInnen}{Herausgegeben von Jahnke, SelmaMüller, Martin Anton}\input{../tex-inputs/latex-leseansicht-abspann}
